\chapter{Mở đầu}
\label{ch:modau}

\begin{muctieu}
Chương này trình bày bối cảnh, vấn đề cần giải quyết, mục tiêu đo lường được,
phạm vi và bố cục của đồ án Vị Việt.
\end{muctieu}

\section{Lý do chọn đề tài}

Ẩm thực Việt Nam có số lượng món lớn, tên gọi và biến thể khác nhau theo vùng
miền. Khi bắt gặp một món ăn qua ảnh, người dùng có thể không biết tên chính
xác để tìm kiếm. Ngược lại, kết quả tìm kiếm trên Internet thường phân tán,
khó phân biệt nội dung có nguồn với nội dung được sao chép hoặc tạo tự động.
Hai khó khăn này tạo ra nhu cầu về một hệ thống vừa nhận diện món qua ảnh,
vừa dẫn người dùng đến công thức có thể truy nguyên.

Nhận diện món ăn là một bài toán phân loại ảnh nhiều lớp. Food-101 cho thấy
đây là bài toán khó do sự khác biệt lớn về cách trình bày, ánh sáng và nền
ảnh, đồng thời nhiều món có hình thức gần giống nhau
\cite{bossard2014food101}. Với món Việt, 30VNFoods cung cấp 25.136 ảnh thuộc
30 lớp và là nguồn phù hợp để đánh giá model trong phạm vi đề tài
\cite{do2021_30vnfoods}. Đồ án chọn cách khai thác mô hình ViT tiền huấn luyện
thay vì huấn luyện mạng mới từ đầu, nhờ đó tập trung nguồn lực vào tích hợp
hệ thống, chất lượng dữ liệu công thức và trải nghiệm mobile.

Sản phẩm được đặt tên \term{Vị Việt}, hướng đến trải nghiệm tối giản: người
dùng có thể tìm món bằng chữ hoặc chụp/chọn ảnh, xem kết quả Top~1 kèm độ tin
cậy, rồi mở đúng công thức tương ứng. Công thức và ảnh chỉ được xuất bản khi
có metadata nguồn và giấy phép hợp lệ.

\section{Mục tiêu của đồ án}

Đồ án hướng tới các mục tiêu cụ thể sau:
\begin{itemize}
  \item Xây dựng ứng dụng Flutter chạy trên Android, hỗ trợ chế độ khách,
        đăng ký, đăng nhập và quản lý phiên.
  \item Tích hợp model ViT-L/16 gồm 38 nhãn món Việt; hỗ trợ ảnh JPEG/PNG,
        trả tên món, độ tin cậy và dữ liệu món được ánh xạ từ database.
  \item Cung cấp danh mục 38 món có ảnh, nguyên liệu, bước nấu, thời gian,
        khẩu phần, nguồn công thức và metadata giấy phép.
  \item Hỗ trợ tìm kiếm, lọc, yêu thích, lịch sử nhận diện, phản hồi nhãn
        đúng và chỉnh sửa tên tài khoản.
  \item Xây dựng backend FastAPI có kiểm tra đầu vào, xác thực JWT, migration,
        audit dữ liệu và tài liệu OpenAPI.
  \item Đánh giá model trên toàn bộ 5.014 ảnh thuộc tập Test sạch của 30 lớp
        giao nhau giữa 30VNFoods và model; báo cáo trung thực giới hạn tám
        nhãn chưa được bộ dữ liệu này bao phủ.
\end{itemize}

\section{Đối tượng sử dụng}

Đối tượng chính là người quan tâm tới món ăn Việt và muốn khám phá công thức
bằng tên món hoặc hình ảnh. Khách có thể xem danh mục, tìm kiếm và nhận diện.
Người dùng đăng nhập có thêm khả năng lưu yêu thích, xem/xóa lịch sử và gửi
phản hồi khi kết quả chưa đúng. Đối tượng kỹ thuật gồm người quản lý dữ liệu
công thức và lập trình viên cần audit, import hoặc triển khai lại hệ thống.

\section{Phạm vi và giới hạn}

Phạm vi nhận diện được giới hạn ở 38 nhãn trong cấu hình model. Giao diện chỉ
hiển thị Top~1 để tối giản; backend vẫn lưu Top~3 cho lịch sử và đánh giá. Ảnh
người dùng chỉ được giữ trong bộ nhớ trong lúc suy luận, không được ghi vào ổ
đĩa hoặc database. Đồ án chưa huấn luyện lại model production, chưa triển khai
backend công khai và chưa tối ưu cho nhiều API instance. SQLite phù hợp với
demo đơn máy nhưng không thay thế hệ quản trị cơ sở dữ liệu cho production.

Dataset 30VNFoods là dữ liệu ảnh phục vụ audit và đánh giá, không phải database
công thức. Tám nhãn model không xuất hiện trong dataset nên các độ đo trên tập
Test chỉ phản ánh 30 lớp giao nhau. Công thức được biên soạn hoặc dịch từ nguồn
có thể truy nguyên, nhưng không được tuyên bố là nội dung đã qua chuyên gia ẩm
thực kiểm duyệt.

\section{Bố cục báo cáo}

Phần còn lại của báo cáo được tổ chức như sau. Chương~\ref{ch:coso} trình
bày cơ sở lý thuyết và công nghệ. Chương~\ref{ch:thietke} phân tích yêu cầu
và thiết kế hệ thống. Chương~\ref{ch:caidat} mô tả quá trình cài đặt, dữ liệu,
giao diện và kiểm thử. Chương~\ref{ch:ketqua} trình bày kết quả đánh giá,
bàn luận và hạn chế. Chương~\ref{ch:ketluan} tổng kết và đề xuất hướng phát
triển.
